% generated from Bio 3/repeat_splitting.tex
\chapter{The sandpile group of a de Bruijn graph splits along its repeats}\label{chap:Bio3}
\chaptermark{Bio3}
\begin{chapterabstract}
\noindent
Two companion papers showed that the sandpile group $\K(G_k)$ of a genome's de Bruijn
graph records repeat length, multiplicity and interleaving --- a single repeat of length
$L$ and multiplicity $r$ contributes $(\Z/r)^{L-k}$, two interleaved repeats add a $\Z/2$,
three can add a $\Z/3$ --- and showed that the correction is not a rank of the chord
diagram. Here we prove the general statement. \emph{Bundle-chain lemma:} in
any balanced multidigraph, a chain $v_1\to\cdots\to v_s$ in which every vertex sends all
$r$ of its arcs to the next can be contracted to a point, and the sandpile group splits
off a direct summand $(\Z/r)^{s-1}$ exactly. The proof is a unimodular change of
variables and takes ten lines; the $r=1$ case is the familiar unitig contraction.
Contracting every bundle chain leaves a \emph{skeleton} with one vertex per repeat family
and one arc per adjacency of repeat occurrences around the genome --- a line-digraph
dual of the repeat graph of Pevzner, Tang and Tesler --- and
\[
\K(G_k)\;\cong\;\bigoplus_{\text{chains }c}(\Z/r_c)^{\,\ell_c-1}\;\oplus\;\K(\Sk).
\]
The skeleton is canonical, and the repeat graph itself is the case in which every chain
has length two. The first summand is what longer $k$-mers remove; the second, the
sandpile group of the repeat-adjacency skeleton, is what no $k$ removes. For repeats of
multiplicity two the second summand is the critical group of the one-vertex map whose
edges are the repeats: it is $\coker(A+I)$ for Bouchet's signed interlace matrix $A$
(Merino--Moffatt--Noble), its order is the interlace polynomial at $1$
(Arratia--Bollob\'as--Sorkin), and --- as their example shows and we recompute --- it is
not a function of the unsigned chord-crossing graph, although it happens to be for up to
five repeats. For $\varphi$X174 the decomposition is
verified against direct computation at $k=10,11,12$ and extends the computable spectrum
to $k=9$, where the compacted graph ($146$ branch vertices) exceeded our exact Smith-form
range but its skeleton does not: $\K(G_9)=(\Z/2)^{42}\oplus(\Z/4)^{4}\oplus\Z/7\oplus
\cdots\oplus\Z/2707831823$, of which $(\Z/2)^{36}$ is length-driven and the rest is the
arrangement of repeats. Seven checks, all exact.
\end{chapterabstract}

\section{Introduction}\label{Bio3:sec:intro}

For a circular sequence $S$ and $k\ge1$, $G_k(S)$ is the de Bruijn multigraph whose
vertices are the distinct $k$-mers and whose arcs are the occurrences of $(k{+}1)$-mers;
$S$ is an Eulerian circuit, and by the BEST theorem \cite{AEdB,SmithTutte} the number of
reconstructions of $S$ from its $k$-mer spectrum is $|\K(G_k)|\prod_v(d^+(v)-1)!$ with
$\K(G_k)$ the sandpile group. The first companion paper \cite{Bio1} showed that $\K(G_k)$
is not determined by the branch statistics and computed it in three model situations: a
single repeat of length $L$ and multiplicity $r$ gives $(\Z/r)^{L-k}$; two repeats of
multiplicity two give $(\Z/2)^{2s-2}$ if they are disjoint or nested and $(\Z/2)^{2s-1}$ if
interleaved, $s=L-k+1$; three repeats give a correction that is $\Z/2$, $\Z/3$ or
$(\Z/2)^{2}$ according to their crossing pattern; since the three non-trivial patterns
have intersection matrices of the same rank over $\F_2$, the correction is not a function
of that rank (Remark~\ref{Bio3:rem:chord} sharpens this). The second
paper \cite{Bio2} used the group as a coordinate system on reconstructions and found that
for $\varphi$X174 at $k=12$ the single surviving $\Z/2$ is a pair of interleaved
$12$-mers.

Those three model computations are special cases of one structural fact, which this note
proves. The fact is local and holds in any balanced multidigraph; its consequence for
genomes is that the sandpile group splits as a direct sum of a \emph{length-driven} part,
one cyclic factor per shared $k$-mer of each repeat, and an \emph{arrangement} part, the
sandpile group of a repeat-adjacency skeleton.

\paragraph{What is known.} Contracting a vertex of in- and out-degree one leaves the
sandpile group unchanged \cite[Lem.~4]{Bio1}; this is the $r=1$ case of
Lemma~\ref{Bio3:lem:bundle} and is the reason assemblers' compacted (unitig) graphs
\cite{BCALM,KM} carry the same group. For \emph{undirected} graphs with multiple edges
the sandpile groups of thick trees \cite{ChenSchedler} and thick cycles \cite{Alar} are
known, and the thick-tree answer, $\bigoplus\Z/m_e$ over the edge multiplicities, looks
like our formula; but the undirected analogue of Lemma~\ref{Bio3:lem:bundle} is false
(Remark~\ref{Bio3:rem:undirected}), so the lemma is not a corollary of those results. The
nearest directed relative is Levine's theorem that $\K(G)$ is a quotient of $\K(LG)$ for
the line digraph \cite{Levine,PSX}. The Eulerian-path formulation of assembly is due to
Pevzner, Tang and Waterman \cite{PTW}, and the \emph{repeat graph}, in which repeat
families are edges of multiplicity $>1$ and vertices are repeat boundaries, to Pevzner,
Tang and Tesler \cite{PTT}; it is the central object of the long-read assembler Flye
\cite{Flye}. Our skeleton has one \emph{vertex} per repeat family and one arc per
adjacency of occurrences, so it stands to the repeat graph roughly as a line digraph to
its graph; we keep the two apart by name. For two-copy repeats the skeleton is the
directed medial graph of a one-vertex map, and its sandpile group is the critical group
of that map in the sense of Merino, Moffatt and Noble \cite{MMN}, which is where the
question of what determines it has already been answered (Remark~\ref{Bio3:rem:chord}). We
have not found the $r>1$ splitting of Lemma~\ref{Bio3:lem:bundle} in the sandpile literature,
nor the sandpile group of the repeat graph considered as such; a search cannot establish
a negative.

\section{Setting}\label{Bio3:sec:setting}

All digraphs are finite multidigraphs, balanced ($d^-(v)=d^+(v)$ at every vertex) and
connected. The out-Laplacian is $\Delta=D^+-A$; the sandpile group is
$\K(G)=\coker(\Delta_{\hat s}^{\mathsf T})$, the row and column of a sink $s$ deleted; for
balanced $G$ the isomorphism class does not depend on $s$ \cite[Lem.~4.12]{HLMPPW}, and
$|\K(G)|$ is the number of spanning arborescences toward $s$. We compare finite abelian
groups by their elementary divisors.

\begin{definition}[Bundle chain]\label{Bio3:def:chain}
A \emph{bundle chain} of multiplicity $r$ and length $s\ge2$ is a sequence of distinct
vertices $v_1,\dots,v_s$ such that for each $t<s$ the vertex $v_t$ has out-degree $r$ and
all $r$ of its out-arcs go to $v_{t+1}$, and $v_s$ has out-degree $r$. Since $G$ is
balanced, every in-arc of $v_{t+1}$ ($t<s$) comes from $v_t$. The \emph{contraction}
$G/C$ replaces the chain by a single vertex carrying the in-arcs of $v_1$ and the
out-arcs of $v_s$; it is again balanced, and an arc $v_s\to v_1$ becomes a loop, which
does not affect $\Delta$.
\end{definition}

\section{The lemma}\label{Bio3:sec:lemma}

\begin{lemma}[Bundle-chain contraction]\label{Bio3:lem:bundle}
Let $C=(v_1,\dots,v_s)$ be a bundle chain of multiplicity $r$ in a balanced multidigraph
$G$ with at least one vertex outside $C$. Then
\[
\K(G)\;\cong\;(\Z/r)^{s-1}\;\oplus\;\K(G/C).
\]
\end{lemma}

\begin{proof}
Root at a vertex outside $C$ and write $x_v$ for the generator of the cokernel
presentation at each non-root vertex, the relations being the rows of $\Delta$. Change
generators on the chain:
\[
z:=x_{v_1},\qquad y_t:=x_{v_t}-x_{v_{t+1}}\quad(1\le t<s),
\]
so that $x_{v_t}=z-\sum_{u<t}y_u$ and in particular $x_{v_s}=z-\sum_{u<s}y_u$; this is a
unimodular (triangular) change of basis. The rows of $\Delta$ fall into three kinds.
\begin{itemize}
\item The row of $v_t$ for $t<s$ is $r\,x_{v_t}-r\,x_{v_{t+1}}=r\,y_t$.
\item The row of $v_s$ is $r\,x_{v_s}-\sum_w c_{v_sw}x_w
=r\,z-r\sum_{u<s}y_u-\sum_w c_{v_sw}x_w$. Subtracting the rows $r\,y_u$ --- a unimodular
row operation --- leaves $r\,z-\sum_w c_{v_sw}x_w$, which is the row of the contracted
vertex in $\Delta(G/C)$.
\item The row of a vertex $u\notin C$ involves $x_{v_t}$ only through arcs from $u$ into
the chain, and all of those land on $v_1$; so it reads $d^+(u)x_u-\sum_w c_{uw}x_w$ with
$x_{v_1}=z$, which is its row in $\Delta(G/C)$.
\end{itemize}
The relation lattice is therefore generated by $\{r\,y_t\}_{t<s}$ together with the rows
of $\Delta(G/C)$ in the variables $z$ and $x_u$, and the generators separate accordingly.
Hence $\K(G)\cong\bigoplus_{t<s}\Z/r\oplus\K(G/C)$.
\end{proof}

If $G$ consists of the chain alone --- the only vertices are $v_1,\dots,v_s$ and the $r$
arcs of $v_s$ return to $v_1$ --- the same computation rooted at $v_1$ gives
$\K(G)\cong(\Z/r)^{s-1}$; this is the single-repeat theorem of \cite[Thm.~5]{Bio1}.

\begin{remark}\label{Bio3:rem:r1}
For $r=1$ the summand is trivial and the lemma says that smoothing a vertex of in- and
out-degree one preserves the group, which is \cite[Lem.~4]{Bio1}. The content of the
lemma for $r>1$ is that a vertex which \emph{cannot} be smoothed --- it has $r$ parallel
exits --- can nevertheless be removed at the exact cost of a $\Z/r$.
\end{remark}

\begin{remark}[The direction is essential]\label{Bio3:rem:undirected}
The undirected statement one might guess --- a path of doubled edges contributes
$(\Z/2)^{\text{length}-1}$ --- is false. The undirected $4$-cycle with every edge doubled
has reduced Laplacian $2M$ with $M$ the reduced Laplacian of the simple $4$-cycle, whose
Smith form is $\operatorname{diag}(1,1,4)$; so its sandpile group is
$(\Z/2)^{2}\oplus\Z/8$, whereas the directed doubled $4$-cycle is a single bundle chain and
has $(\Z/2)^{3}$. What makes the directed lemma work is that every arc \emph{into} the
chain lands on its head $v_1$ and every arc \emph{out} of it leaves from $v_s$ with the
full multiplicity $r$; an undirected doubled edge carries flow both ways and the change
of variables in the proof no longer separates.
\end{remark}

\section{The splitting theorem}\label{Bio3:sec:split}

Contracting bundle chains one after another terminates, since each contraction reduces
the number of vertices, and the result contains no bundle chain of length $\ge2$. We call
it the \emph{skeleton} $\Sk(G)$. Its sandpile group is well defined by the lemma; the
skeleton itself is too.

\begin{lemma}[The skeleton is canonical]\label{Bio3:lem:canon}
Call $v\to w$ a \emph{chain step} if $v\ne w$, all out-arcs of $v$ go to $w$, and
$d^+(w)=d^+(v)$. In a balanced multidigraph the chain steps form a partial injection on
the vertex set, so the maximal bundle chains are its maximal paths: they are pairwise
vertex-disjoint and determined by $G$. Contracting them in any order gives the same
digraph up to isomorphism, which is chain-free.
\end{lemma}

\begin{proof}
If $v\to w$ is a chain step then $w$ receives $d^+(v)=d^+(w)=d^-(w)$ arcs from $v$, hence
all of its in-arcs, so no $v'\ne v$ steps to $w$; and $v$ has a single target. Thus the
step relation is a partial injection and its maximal paths partition the vertices it
touches (a path may close into a cycle only if $G$ is a single closed chain, when the
skeleton is one vertex with $r$ loops). Contracting a maximal chain $C$ to a vertex $c$
changes no degree outside $C$, gives $c$ the in-arcs of $v_1$ and the out-arcs of $v_s$,
and so creates no new chain step: $u\to c$ would require $u\to v_1$ to have been a step
and $c\to w$ would require $v_s\to w$ to have been one, both excluded by maximality. Nor
does it destroy a step outside $C$. Hence the contractions of distinct maximal chains
commute, and the result is $G$ with every maximal chain contracted.
\end{proof}

Check~(E) verifies this on $100$ random balanced digraphs with planted chains: the
partition into maximal chains, the multiset of $(r_c,\ell_c)$ and the skeleton's group are
invariant under random relabelling, and the skeleton contains no chain step.

\begin{theorem}[Splitting]\label{Bio3:thm:split}
For any balanced connected multidigraph $G$, with $\{(r_c,\ell_c)\}$ the multiplicities
and lengths of the chains contracted on the way to a skeleton,
\[
\K(G)\;\cong\;\bigoplus_c(\Z/r_c)^{\ell_c-1}\;\oplus\;\K(\Sk(G)),
\qquad
|\K(G)|=\prod_c r_c^{\,\ell_c-1}\cdot|\K(\Sk(G))|.
\]
\end{theorem}

\begin{proof}
Induction on the number of contractions, using Lemma~\ref{Bio3:lem:bundle} at each step.
\end{proof}

Now let $S$ be a sequence with repeats $W_1,\dots,W_m$ of lengths $L_i$ and
multiplicities $r_i$, and suppose the \emph{clean} hypothesis of \cite{Bio1}: the $k$-mers
inside $W_i$ occur exactly $r_i$ times each and every other $k$-mer once. Write $w$ for
the circular word of repeat occurrences and $D_w$ for the \emph{connector digraph}: one
vertex per repeat, one arc $i\to j$ for each occurrence of $W_i$ immediately followed
around the circle by an occurrence of $W_j$ (loops deleted). $D_w$ is balanced with
out-degree $r_i$ at vertex $i$.

\begin{corollary}[Clean sequences]\label{Bio3:cor:clean}
Under the clean hypothesis, with $s_i=L_i-k+1$,
\[
\K\bigl(G_k(S)\bigr)\;\cong\;\bigoplus_{i=1}^{m}(\Z/r_i)^{\,s_i-1}\;\oplus\;\K(D_w).
\]
\end{corollary}

\begin{proof}
After unitig contraction the graph consists of the $m$ spines --- bundle chains of
multiplicity $r_i$ and length $s_i$ --- joined by the connector arcs; contracting the
spines gives $D_w$. Apply Theorem~\ref{Bio3:thm:split}.
\end{proof}

\begin{corollary}[The published cases]\label{Bio3:cor:cases}
For $m=1$, $D_w$ is a single vertex and $\K=(\Z/r)^{s-1}$ \cite[Thm.~5]{Bio1}. For $m=2$,
$r_1=r_2=2$: $D_w$ is trivial for $A\,A\,B\,B$ and $A\,B\,B\,A$ and is the doubled
$2$-cycle, with $\K(D_w)=\Z/2$, for $A\,B\,A\,B$ \cite[Thm.~6]{Bio1}. For $m=3$ the four
crossing patterns give $\K(D_w)=0,\ \Z/2,\ \Z/3,\ (\Z/2)^{2}$ \cite[Prop.~8]{Bio1}; the
$\Z/3$ is the sandpile group of the complete symmetric digraph on three vertices, which is
$D_w$ for $A\,B\,A\,C\,B\,C$.
\end{corollary}

\begin{remark}[The skeleton and the repeat graph]\label{Bio3:rem:rg}
The duality with the repeat graph of \cite{PTT,Flye} can be made exact. Under the clean
hypothesis, collapse each spine of the compacted graph to a single arc: the result is the
multidigraph $\mathrm{RG}_w$ with two vertices $h_i,t_i$ per family, $r_i$ parallel arcs
$h_i\to t_i$, and one arc $t_i\to h_j$ per adjacency $i\,j$ in $w$ --- repeats as
multi-arcs between their boundary vertices, unique segments as single arcs, which is the
repeat graph. Each pair $(h_i,t_i)$ is a bundle chain of length two, and contracting it
identifies the multi-arc with a vertex: $D_w$ is the transition digraph of the repeat
arcs of $\mathrm{RG}_w$ (the line digraph, restricted to repeat arcs, with unique
segments as the adjacencies). By Lemma~\ref{Bio3:lem:bundle},
\[
\K(\mathrm{RG}_w)\;\cong\;\bigoplus_{i=1}^{m}\Z/r_i\;\oplus\;\K(D_w),
\]
so the sandpile group of the repeat graph is the arrangement part plus one $\Z/r_i$ per
family --- the length-driven part at $s_i=2$ (check~(E), $200$ random words). The two
objects carry the same information; we work with $D_w$ because it has no length-driven
residue.
\end{remark}

\begin{remark}[Which invariant of the chord diagram]\label{Bio3:rem:chord}
For $r_i=2$ the $2m$ occurrences form a chord diagram with $m$ chords --- a
double-occurrence word, or a one-vertex map (bouquet) with the repeats as edges --- and
$D_w$ is a $2$-in $2$-out digraph. On what invariant of the diagram does $\K(D_w)$
depend? Not on the rank of the unsigned intersection matrix, over $\F_2$ or over $\Z$,
and not on its integer Smith form: for $m=3$ the three non-trivial patterns all have
$\F_2$-rank $2$, and $A\,A\,B\,C\,B\,C$ (one crossing) and $A\,B\,A\,C\,B\,C$ (path)
share the Smith form $(1,1,0)$ while their groups are $\Z/2$ and $\Z/3$; for $m=5$ the
rank-$2$ diagrams carry nine different groups, among them $\Z/7$ and $\Z/9$
(check~(D)). The \emph{order} $|\K(D_w)|$ is the number of Eulerian circuits of $D_w$,
and this is a function of the intersection graph $H$ alone (the \emph{interlace graph}
of \cite{ABS}, whose Definition~1 is exactly chord crossing): it is the interlace
polynomial $q(H;1)$ \cite[Thm.~9, Cor.~13]{ABS}, with roots in Martin's and Las Vergnas'
circuit-partition polynomials \cite{LasVergnas}. Exhaustively for $m\le5$ the group is
also constant on the $4$, $11$ and $34$ isomorphism classes of $H$ --- but this is an
accident of small $m$. Merino, Moffatt and Noble \cite[Rem.~6.7]{MMN} give two
$6$-chord bouquets, $1\,2\,3\,4\,5\,6\,1\,3\,2\,4\,6\,5$ and
$1\,2\,3\,6\,5\,4\,1\,3\,2\,5\,6\,4$, with the same interlace graph ($K_6$ minus two
disjoint edges, so $|\K|=18$ for both) and critical groups $\Z/3\oplus\Z/6$ and
$\Z/18$; check~(D) recomputes both from $D_w$ and exhibits two further such pairs. What
does determine the group is the \emph{signed} interlace matrix $A$ of Bouchet
($A_{ef}=\pm1$ according as the occurrences run $e\,f\,e\,f$ or $f\,e\,f\,e$, $0$ if
not interlaced): by \cite[Thm.~7.3]{MMN}, whose directed medial graph of a bouquet is
our $D_w$,
\[
\K(D_w)\;\cong\;\coker(A+I),
\]
the critical group of the map; its order $\det(A+I)$ as an Euler-circuit count is older
\cite{MacrisPule,Jonsson}. Check~(D) confirms the isomorphism on every double-occurrence
word with $m\le5$ and on $1200$ random words with $6\le m\le9$. So the arrangement part
of an assembly's ambiguity, for two-copy repeats, is the critical group of the bouquet
whose edges are the repeats, computable from a signed $m\times m$ matrix rather than
from the de Bruijn graph; and the reason it is not a function of the chord-crossing
graph is the reason \cite{MMN} record: the signs matter. For $r_i>2$ the occurrences
form a hypergraph rather than a chord diagram and we know of no analogue.
\end{remark}

\begin{remark}[What $k$ can and cannot remove]\label{Bio3:rem:k}
Increasing $k$ shortens every spine: $s_i$ decreases by one per unit of $k$, and the
length-driven summand $\bigoplus(\Z/r_i)^{s_i-1}$ shrinks accordingly, vanishing at
$k=\max L_i$. The skeleton summand $\K(D_w)$ does not depend on $k$ at all, as long as
$k\le\min_iL_i$. So the global ambiguity of an assembly splits into a part that longer
$k$-mers resolve and a part that only information beyond the repeats --- long reads,
pairs, a reference --- can resolve, and the second part is exactly the sandpile group of
the repeat-adjacency skeleton. By BEST applied to $D_w$, $|\K(D_w)|\cdot\prod_i(r_i-1)!$ is the number
of Eulerian circuits of $D_w$, i.e.\ of circular words with the same adjacency counts as
$w$: the number of ways to thread the repeats, ignoring their lengths. That this is
$|\K(D_w)|$ up to the local factor is the repeat-level form of the decomposition of
\cite{Bio1}.
\end{remark}

\section{$\varphi$X174}\label{Bio3:sec:phix}

Table~\ref{Bio3:tab:phix} applies the theorem to the first sequenced genome (NC\_001422.1,
$5386$ bp). At each $k$ the compacted graph is bundle-contracted; the chains, the
skeleton, and the two summands are listed, and where the compacted graph is small enough
the group is also computed directly for comparison. Two things happen. The direct and
decomposed groups agree wherever both exist ($k=10,11,12$). And at $k=9$, where the
compacted graph has $146$ branch vertices and was beyond the dense exact Smith form used
in \cite{Bio1}, the skeleton has $110$ and is not: the spectrum extends by one level, with
the length-driven part $(\Z/2)^{36}$ split off exactly and the arrangement part read off
the skeleton.

\begin{table}[ht]
\centering\footnotesize
\begin{tabular}{rrrrll}
\toprule
$k$ & branch & chains & skeleton & $\K(\Sk)$ (arrangement part) & length-driven part\\
\midrule
12 &   2 &  1 &   1 & $0$ & $\Z/2$\\
11 &  10 &  2 &   8 & $\Z/3\oplus\Z/11$ & $(\Z/2)^{2}$\\
10 &  47 &  8 &  37 & $\Z/4\oplus\Z/1076338579$ & $(\Z/2)^{10}$\\
 9 & 146 & 28 & 110 & $(\Z/2)^{6}\oplus(\Z/4)^{4}\oplus\Z/7\oplus\Z/523\oplus\Z/3779$ & $(\Z/2)^{36}$\\
   &     &    &     & $\quad\oplus\Z/5297\oplus\Z/6561\oplus\Z/7669\oplus\Z/2707831823$ & \\
 8 & 401 & 65 & 312 & (skeleton beyond the exact-SNF range used) & \\
\bottomrule
\end{tabular}
\caption{$\varphi$X174 under bundle-chain contraction: number of branch vertices of the
compacted graph, number of bundle chains contracted, number of skeleton vertices, and the
two summands. $\K(G_k)$ is their direct sum; at $k=10,11,12$ this agrees with the group
computed directly on the compacted graph (check~(F)). All chains have multiplicity $2$.
All entries exact.}
\label{Bio3:tab:phix}
\end{table}

The reading at $k=11$ is the cleanest: $\K(G_{11})=(\Z/2)^{2}\oplus\Z/3\oplus\Z/11$, and
the theorem says which is which --- the two $\Z/2$'s are two repeat families each with one
extra shared $11$-mer beyond the minimum, and $\Z/3\oplus\Z/11$, of order $33$, is the
arrangement ambiguity of the eight-vertex skeleton, which no choice of $k\le11$ can
reduce.

\section{Verification}\label{Bio3:sec:battery}

The script \texttt{repeat\_splitting.py} runs the seven checks of Table~\ref{Bio3:tab:battery}
in about thirty seconds, all passing; its output is archived as
\texttt{repeat\_splitting\_output.txt}. All arithmetic is exact; groups are compared by
elementary divisors.

\begin{table}[ht]
\centering\small
\begin{tabular}{clc}
\toprule
key & check & mode\\
\midrule
(A) & Corollary~\ref{Bio3:cor:clean} on $120$ clean random sequences, $m\le4$, $r_i\le4$ & exact\\
(B) & Lemma~\ref{Bio3:lem:bundle} on $60$ random balanced digraphs with a planted chain, $r\le6$, $s\le4$ & exact\\
(C) & Corollary~\ref{Bio3:cor:cases}: the nine published values are the formula on $D_w$ & exact\\
(D) & Remark~\ref{Bio3:rem:chord}: all double-occurrence words, $m\le5$; rank, Smith form, graph & exact\\
(E) & Lemma~\ref{Bio3:lem:canon} on $100$ relabelled digraphs; Remark~\ref{Bio3:rem:rg} on $200$ words & exact\\
(F) & Table~\ref{Bio3:tab:phix}: decomposition vs.\ direct computation; extension to $k=9$ & exact\\
\bottomrule
\end{tabular}
\caption{The verification battery, $7/7$ (check~(E) has two rows).}
\label{Bio3:tab:battery}
\end{table}

\section{Scope}\label{Bio3:sec:scope}

\emph{Proved.} Lemma~\ref{Bio3:lem:bundle}, Theorem~\ref{Bio3:thm:split} and its corollaries, for
arbitrary balanced multidigraphs and arbitrary multiplicities; the clean hypothesis
enters only in identifying the skeleton with $D_w$.

\emph{Quoted.} Sink-independence \cite{HLMPPW}; BEST \cite{AEdB,SmithTutte}; the
Eulerian formulation and the repeat graph \cite{PTW,PTT,Flye}; the thick-graph and
line-digraph results that are the lemma's nearest relatives
\cite{ChenSchedler,Alar,Levine,PSX}; the special cases and the $\varphi$X174 spectrum
\cite{Bio1,Bio2}.

\emph{Open.} (a) A structure theorem for $\K(D_w)$ itself --- for which circular words is
it trivial, and what bounds its exponent? For $r_i=2$ this is the critical group of a
bouquet (Remark~\ref{Bio3:rem:chord}), and \cite{MMN} note that for $3$-connected regular
delta-matroids the group is well defined on the delta-matroid; whether $\K(D_w)$ is
determined by the interlace graph whenever that graph is prime (no split) we have not
examined. For $r_i>2$ nothing of the kind is available. (b) Real genomes are not clean: the
skeleton of $\varphi$X174 at $k=9$ has $110$ vertices and carries factors as large as
$2707831823$. All six factors above $100$ in Table~\ref{Bio3:tab:phix} ($523$, $3779$, $5297$,
$7669$, $1076338579$, $2707831823$) are prime, which is what the sandpile group of a
generic digraph looks like (typically cyclic, or nearly so); we see no reason to expect a
topological reading of a large prime cyclic factor and claim none. (c) Double strands: reverse complementation reverses arc
direction, so a bidirected version of Lemma~\ref{Bio3:lem:bundle} must start from the signed
Laplacian; we have not attempted it.

\section*{Provenance of this chapter}

Unrefereed preprint. Lemmas~\ref{Bio3:lem:bundle} and~\ref{Bio3:lem:canon} are proved in full
above and are the only new mathematics; Remark~\ref{Bio3:rem:chord} is a translation of
published results \cite{ABS,MMN} into the assembly setting, verified computationally ---
an earlier draft of this note conjectured that the group was an invariant of the unsigned
interlace graph, which \cite[Rem.~6.7]{MMN} had already shown to be false; everything
else is either quoted or an exact computation reproducible in seconds. We searched for the $r>1$ splitting in the sandpile literature and for any use of
the sandpile group of a repeat graph or repeat-adjacency skeleton and found neither; the
undirected thick-graph results \cite{ChenSchedler,Alar} are the closest, and
Remark~\ref{Bio3:rem:undirected} shows they do not imply the lemma.
